\chapter{Bread for tomorrow}\label{connection:keeping}
\regionline{Drying, carrying and returning to the table}
\atlasfiguretitle{Bread for tomorrow: paired comparison}
\comparisonplate{Carasau}{Knack}{Two crisp breads. Puffing and splitting make one stack, rolling makes the other.}

A crisp sheet asks a different question from a soft loaf. Instead of wondering only how bread can be kept fresh, we can ask what kind of bread is useful when it will not be eaten immediately. Pane carasau makes that question visible. A risen round puffs, is split into sheets and returns to the oven to become dry. Its making prepares a different future for the bread.\citep{carasaurecipe,carasauculture}

That future is part of the Sardinian bread story, but the work also belongs to a group. SardegnaCultura describes coordinated roles among women, with relatives or neighbours helping and assistance repaid or rewarded with bread. The finished stack therefore holds more than a clever response to storage. Making enough bread can require a household to draw other people into the bake.\citep{carasauculture}

\section{Making bread wait}
Knäckebröd brings a northern comparison into view. The thin, dried bread in the Swedish account could hang on rods. Its recipe in this book is rolled and baked as crispbread, rather than puffed and split like carasau. Dryness links their use, while the dough and handling keep the making distinct. A shared purpose can be reached through different sequences.\citep{duxknack,daumknack}

The Finnish reikäleipä story supplies another kind of stored bread. Here a flattened rye form and a central opening make sense alongside the account of pole storage. The rye loaf recipe printed in the atlas is a round, so readers should return to the story for that variety rather than expecting the recipe to demonstrate every form named there.\citep{finroot}

Keeping bread changes its place in the home. A stack of sheets, a hanging round and sliced bread in a freezer need different spaces. Each also changes the moment at which work is done. A larger bake can move effort to one day, while small fresh breads spread it across several meals. Neither arrangement is effortless. The useful comparison is between rhythms of work and eating.

\section{A dry bread can come back differently}
Lavash provides a particularly appealing turn. UNESCO's Armenian account describes dried bread being softened with water for use again. Its thinness allows keeping and later flexibility to belong to the same bread. A sheet need not be discarded simply because it has lost the softness it had near the oven.\citep{lavash}

Pane frattau takes carasau into a meal with broth, tomato, pecorino and egg in the Sardinian account. Moisture changes the sheet's task. The bread made to snap becomes a layer in a dish, holding flavour and giving the meal substance. A crisp bread does not always need to remain crisp to be used well.\citep{sardiniabreads}

Casabe gives the atlas a root-based comparison. The cassava-bread chapter follows prepared dry bread into the home kitchen, where its serving recipe uses garlic and oil. The history includes skilled processing before the bread reaches the eater. Its crispness is a useful point of comparison with carasau, but it does not make raw cassava processing interchangeable with wheat dough preparation. Buy finished casabe for the recipe here and keep the distinction visible.\citep{cassava,claracasabe08}

\section{What travels can be unfinished}
A sack of flour and a baked loaf have different uses on a journey. The damper story includes wheat bread made around a campfire, where the cook turns carried ingredients into food at the place of eating. In that account, travel changes the location of breadmaking rather than requiring every loaf to be baked in advance.\citep{worldchefsdamper09}

The same chapter reaches farther back to Aboriginal seed breads and the labour of gathering, separating and grinding native seeds. Australian Museum writer Stan Florek makes that chain of work visible. A quick final cooking can rest on extensive knowledge and preparation. It would be a mistake to see a short baking time and imagine an uncomplicated food system.\citep{florekbread09}

A starter can travel too. Bergman's dried rye dough belongs in this comparison because it carries the means of making future bread. The traveller brings a small preparation, then uses the new kitchen's flour and water. In contrast, a purchased crispbread arrives as a finished food. Both can link homes across distance, but they ask different things of the person carrying them.\citep{finjourney}

\section{Tomorrow belongs to the recipe}
For the home baker, the storage note is part of choosing a bake. Will the bread be shared while warm, frozen in useful portions or kept crisp in an airtight container? Is it already filled with cheese or another perishable ingredient? Those questions help determine how much to make and when to serve it. The kitchen notes on page~197 give the practical route through the atlas.

Try a crisp sheet in two meals before deciding that one batch has become repetitive. Its first use may reward a snap. Its next may reward the way it takes up moisture. The bread can remain recognisable while its place on the plate changes, and that is an especially satisfying form of continuity.

\section{A thicker loaf has another rhythm}
The Danish rural rye bread described by Lex was made as stiff dough for loaves that needed to last between widely spaced baking days. Their squared sides came from loaves set beside one another, rather than from separate tins. That account makes keeping visible in the size of the bake and the arrangement inside the oven.\citep{lexrye}

Lex also describes rural bakeries taking on heavy mixing around the mills. A household's bread rhythm could change when somebody else did the work. Keeping, buying and sharing labour belong in the same discussion. A familiar loaf on the table can continue while the task of producing it moves into another pair of hands.

\begin{breadnote}{A sheet in two meals}
Bake carasau on page~29. Keep a portion dry for serving with cheese, then follow the story's pane frattau invitation for a softer meal. For another grain and another drying method, try knäckebröd on page~48.
\end{breadnote}
